\begin{ledgroup}
\footnotesize
\pstart
\noindent%
Die vorliegenden Stücke N.\;\ref{60631} und N.\;\ref{41823} ("Solutio problematis R.\,P. Vannii S.\,J.") bezeugen eine weitere Phase von Leibnizens Beschäftigung 
mit den kontroversen Thesen des Jesuiten
\indexN{\textls{Vanni}, Giovanni Francesco 1638\textendash1709}Giovanni Francesco Vannis
und der von ihm aufgeworfenen Frage der korrekten Zerlegung des Gewichts einer Kugel auf zwei schiefen Ebenen (\textsc{G.\,F.\;Vanni}, \cite{01995}\textit{Specimen libri de momentis gravium}, Rom 1684) um 1690. 
Die zwei Konzepte, die in N.\;\ref{41823} als Zeugen kollationiert werden (der Zettel mit der Signatur LH~XXXV~10, 13~Bl.~9 als \textit{L\textsuperscript{1}}, der Bogen LH~XXXV~10, 13~Bl.~10\textendash11 als \textit{L\textsuperscript{2}}), sind laut Leibnizens eigener Angabe während seiner Rückreise nach Deutschland nach dem \indexO{Italien}italienischen Jahr entstanden ("Dum in \indexO{Deutschland}Germaniam redeo").
Die kurze Aufzeichnung N.\;\ref{60631} enthält keine eindeutigen Datierungselemente, aber sie erweist sich im Vergleich mit den Konzepten als die Vorlage für eine bestimmte Stelle der "Solutio"
(\refpassage{35_10_13_010-011_15a}{35_10_13_010-011_15b}).
\pend
\pstart
\edlabel{intro_vanni-solutio-problematis_1a}%
Nachdem er während seines kurzen Besuchs in \indexO{Parma}Parma am 2./3. Februar 1690
\indexN{\textls{Bacchini}, Benedetto 1651\textendash1721}Benedetto Bacchini 
kennengelernt hatte (siehe den Brief vom 23. März 1690, \cite{02182}\textit{LSB} I,~5 N.\;316, S.\;554\,f.),
plante Leibniz eine Veröffentlichung in der
von Bacchini und \indexN{\textls{Roberti}, Gaudenzio 1655\textendash1695}Gaudenzio Roberti
herausgegebenen, kurzlebigen Zeitschrift \cite{02102}\textit{Giornale de' Letterati di Parma} (1686\textendash1690).
Mit der Vermittlung \indexN{\textls{Magliabechi} (Magliabecchi, Magliabecki, Magliabekius, Magliabequius), Antonio 1633\textendash1714}Magliabechis -- der als Gönner \indexN{\textls{Vanni}, Giovanni Francesco 1638\textendash1709}Vannis um 1684 dem \cite{}\textit{Specimen} zu größerer Verbreitung geholfen hatte (siehe N.\;\ref{intro_vanni-specimen}) -- sollte die "Solutio ... ad clarissimos editores \textit{Diarii Eruditorum} Parmensis ... missa" in der italienischen Zeitschrift erscheinen.
Gemäß einem später verworfenen Konzept von Mitte Mai zu einem Brief an Magliabechi (die abgefertigte Fassung ist vom \cite{02162}16. Juli 1690, \textit{LSB} I,~5 N.\;362, S.\;619--621) hoffte Leibniz,
dass sein Beitrag der "controversia a R. P. Vannio mota", die das italienische und europäische Fachpublikum rege ausgetragen hatte,
ein Ende setzen könnte
(ebd. S.\;621).
\pend
\pstart
\edlabel{intro_vanni-solutio-problematis_1b}
Allerdings sollte die "Solutio" nie in der geplanten Form erscheinen.
Grund dafür ist nicht lediglich, dass das \cite{02102}\textit{Giornale de' Letterati} im Laufe des Jahres 1690 eingestellt wurde, als  \indexN{\textls{Bacchini}, Benedetto 1651\textendash1721}Bacchini \indexO{Parma}Parma verlassen musste (in den Jahren 1692--1698 gab er in \indexO{Modena}Modena ein neues \cite{02103}\textit{Giornale de' Letterati} heraus).
Vielmehr sprachen Leibnizens Bedenken gegen die Publikation.
Zwar hatte er im ersten Teil des Konzepts \textit{L\textsuperscript{2}} (\refpassage{35_10_13_010-011_16b}{35_10_13_010-011_16b}) den Text des ersten Zeugen \textit{L\textsuperscript{1}} eingearbeitet, seinen Aufsatz sorgfältig ausgearbeitet und mit akkuraten Diagrammen (\textit{fig. 1} -- \textit{fig. 3}) versehen, so dass dieser einen formal und inhaltlich sauberen, gepflegten und durchaus druckreifen Text bildet.
Aber im übrigen Teil des Konzepts, gleich im Anschluss an das Ende des Beweises, hat Leibniz neue Schwierigkeiten erblickt, die ihn dazu bewogen, nachdem der Text zu keinem befriedigenden Abschluss führte, diesen einzubehalten und von seinem Veröffentlichungsplan abzusehen.
So schreibt er im Anschluss an den Beweis, auf \refpassage{35_10_13_010-011_12b}{35_10_13_010-011_12b}: "Re accurate expensa video insignem nasci difficultatem" und stellt an mehreren Orten neue Fehler und Schwierigkeiten fest (\refpassage{35_10_13_010-011_13a}{35_10_13_010-011_13b}, \refpassage{35_10_13_010-011_error-specialis_a}{35_10_13_010-011_error-specialis_b} und  \refnote{FnA:35_10_13_010-011_hoc-falsum}).
\pend
\pstart
Zur genauen Datierung der Entstehung des Konzepts \textit{L\textsuperscript{2}} von N.\;\ref{41823} lässt sich auf Grund der Angabe "Dum in Germaniam redeo" und v.\,a. des obengenannten, nicht versendeten Briefkonzept für Magliabechi von Mitte Mai (\textit{LSB} I,~5 N.\;362), Folgendes feststellen:
Leibniz hatte seit frühestens Anfang April am Aufsatz gearbeitet, und ihn bis Mitte Mai mit sehr hoher Wahrscheinlichkeit abgeschlossen (d.\,h. den ersten Teil des Konzepts, \refpassage{35_10_13_010-011_16b}{35_10_13_010-011_16b}) und bereitete die Veröffentlichung vor. 
Im Briefkonzept für Magliabechi, dem er die "Solutio" beilegen wollte, bat Leibniz, den Aufsatz zuerst \indexN{\textls{Bodenhausen} (Bodenus, Bodenausen) Rudolf Christian, Freiherr v. 1698}Bodenhausen mitzuteilen, wohl damit dieser eine Reinschrift für den Druck anfertigen könne (siehe die editorische Erläuterung zu Leibnizens \cite{02163}Brief an Bodenhausen von Ende Juni 1690, \textit{LSB} III~4 N.\;263). Auf diese Weise gingen zur selben Zeit Leibniz und Bodenhausen mit den Entwürfen zur \cite{01354}\textit{Dynamica de potentia} vor.
Aus dem Briefwechsel geht kein konkreter Endpunkt der Arbeit an N.\;\ref{41823} hervor; womöglich war Leibniz schon Mitte Mai auf die neuen Schwierigkeiten gestoßen, vielleicht erst kurze Zeit danach.
Die vorgeschlagene Datierung trägt dem Umstand Rechnung, dass Leibniz in seinem aus Hannover tatsächlich abgeschickten \cite{02162}Brief an Magliabechi seinen frühere Absicht, die "Solutio" zu veröffentlichen, nicht mehr erwähnt.
\pend
\pstart
Für die Aufzeichnung N.\;\ref{60631}, die als Vorlage für eine Passage von N.\;\ref{41823} diente, erscheint die Annahme angemessen, dass diese auch im unmittelbaren Vorfeld des Konzepts entstanden sei, allerdings vor dem Zeitpunkt, zu dem dieses abgeschlossen vorlag (Mitte Mai).
\pend
\end{ledgroup}
